% =====================================================================
%  preamble.tex —— 导言区（Lead 独占；译者只读，不得修改）
%  Mark Gross, "Tropical Geometry and Mirror Symmetry"
%  (AMS CBMS Regional Conference Series in Mathematics 114, 2011)
%  中文译注本 · LuaLaTeX 编译
% =====================================================================
\documentclass[11pt,a4paper]{article}

\usepackage[UTF8]{ctex}
\usepackage{amsmath,amssymb,amsthm,mathtools}
\usepackage{bm}
\usepackage{graphicx}
\usepackage[a4paper,left=2.4cm,right=2.4cm,top=2.6cm,bottom=2.7cm]{geometry}
\usepackage[dvipsnames]{xcolor}
\usepackage{enumitem}
\usepackage{booktabs}
\usepackage{array}
\usepackage{longtable}
\usepackage{multicol}
\usepackage{fancyhdr}
\usepackage{caption}
\usepackage{tcolorbox}
\usepackage{tikz}
\usepackage{titlesec}
\usepackage[hidelinks,unicode]{hyperref}
\hypersetup{pdftitle={热带几何与镜像对称（中文译注本）},
  pdfauthor={Mark Gross 著；中文译注：本工程}}

\captionsetup{font=small,labelsep=quad}
\setlength{\emergencystretch}{3em}
\allowdisplaybreaks[1]
\graphicspath{{figures/}{../figures/}}
% 基域记号：原文用黑板体 \mathbb{k}；为与全文既有译法统一，统一写作 k（斜体）。

% ---------- 篇 / 章 / 节标题 ----------
\renewcommand{\partname}{}
\renewcommand{\thepart}{}
\titleformat{\part}[display]{\centering\Huge\bfseries}{}{0pt}{}
\titlespacing*{\part}{0pt}{40pt}{40pt}

\titleformat{\section}{\Large\bfseries}{第 \thesection{} 章}{1em}{}
\titlespacing*{\section}{0pt}{24pt}{14pt}
\titleformat{\subsection}{\large\bfseries}{\thesubsection}{1em}{}
\titlespacing*{\subsection}{0pt}{16pt}{8pt}
% 原文的三级标题是“行内粗体小标题”（2.1.1. 标题。正文…），用 runin 样式
\titleformat{\subsubsection}[runin]{\bfseries}{\thesubsubsection.}{0.5em}{}
\titlespacing*{\subsubsection}{0pt}{0pt}{0pt}

% ---------- 页眉页脚 ----------
\pagestyle{fancy}
\fancyhf{}
\fancyhead[L]{\small\itshape Tropical Geometry and Mirror Symmetry}
\fancyhead[R]{\small 热带几何与镜像对称\quad\thepage}
\renewcommand{\headrulewidth}{0.4pt}

% ---------- 定理类环境：全书共享一个计数器，按章（section）编号 ----------
%  原文：Definition 1.1 / Examples 1.2 / Lemma 1.13 / Proposition 1.15 …
%  即所有 “定义/例/引理/命题/定理/推论” 共用一个章节内计数器。
\theoremstyle{plain}
\newtheorem{theorem}{定理}[section]
\newtheorem{proposition}[theorem]{命题}
\newtheorem{lemma}[theorem]{引理}
\newtheorem{corollary}[theorem]{推论}
\newtheorem{statement}[theorem]{陈述}
\newtheorem{assumption}[theorem]{假设}
\newtheorem{claim}[theorem]{断言}
\newtheorem{conjecture}[theorem]{猜想}
\newtheorem{question}[theorem]{问题}
\theoremstyle{definition}
\newtheorem{definition}[theorem]{定义}
\newtheorem{example}[theorem]{例}
\newtheorem{examples}[theorem]{例}
\newtheorem{construction}[theorem]{构造}
\newtheorem{remark}[theorem]{注}
\newtheorem{notation}[theorem]{记号}
\newtheorem{exercise}{习题}

% 行内粗体小标题（原文 "**Some title.**"）→ \para{...}
\newcommand{\para}[1]{\par\medskip\noindent\textbf{#1}\quad\ignorespaces}

% 图注/表注（原文图号是“章内连续编号”，如 Figure 2，手工给出）
\newcommand{\figtitle}[2]{\caption*{\small\textbf{图#1：}#2}}
\newcommand{\tabtitle}[2]{\caption*{\small\textbf{表#1：}#2}}

% 参考文献条目
\newcommand{\refitem}[2]{\par\hangindent=1.5em\hangafter=1\noindent [#1]~#2}

% ---------- 常用宏（只准用这里定义过的，不要自造 \newcommand） ----------
\newcommand{\R}{\mathbb{R}}
\newcommand{\C}{\mathbb{C}}
\newcommand{\Z}{\mathbb{Z}}
\newcommand{\N}{\mathbb{N}}
\newcommand{\Q}{\mathbb{Q}}
\newcommand{\PP}{\mathbb{P}}
\newcommand{\cA}{\mathcal{A}}
\newcommand{\cB}{\mathcal{B}}
\newcommand{\cC}{\mathcal{C}}
\newcommand{\cD}{\mathcal{D}}
\newcommand{\cE}{\mathcal{E}}
\newcommand{\cF}{\mathcal{F}}
\newcommand{\cG}{\mathcal{G}}
\newcommand{\cH}{\mathcal{H}}
\newcommand{\cI}{\mathcal{I}}
\newcommand{\cJ}{\mathcal{J}}
\newcommand{\cK}{\mathcal{K}}
\newcommand{\cL}{\mathcal{L}}
\newcommand{\cM}{\mathcal{M}}
\newcommand{\cN}{\mathcal{N}}
\newcommand{\cO}{\mathcal{O}}
\newcommand{\cP}{\mathcal{P}}
\newcommand{\cQ}{\mathcal{Q}}
\newcommand{\cR}{\mathcal{R}}
\newcommand{\cS}{\mathcal{S}}
\newcommand{\cT}{\mathcal{T}}
\newcommand{\cU}{\mathcal{U}}
\newcommand{\cV}{\mathcal{V}}
\newcommand{\cW}{\mathcal{W}}
\newcommand{\cX}{\mathcal{X}}
\newcommand{\cY}{\mathcal{Y}}
\newcommand{\cZ}{\mathcal{Z}}
\newcommand{\fg}{\mathfrak{g}}
\newcommand{\fh}{\mathfrak{h}}
\newcommand{\fM}{\mathfrak{M}}
\newcommand{\fN}{\mathfrak{N}}
\newcommand{\fT}{\mathfrak{T}}
\newcommand{\fX}{\mathfrak{X}}
\newcommand{\fY}{\mathfrak{Y}}
\newcommand{\fZ}{\mathfrak{Z}}
\newcommand{\dd}{\mathrm{d}}
\newcommand{\ee}{\mathrm{e}}
\newcommand{\ii}{\mathrm{i}}
\newcommand{\diag}{\operatorname{diag}}
\newcommand{\sgn}{\operatorname{sgn}}
\newcommand{\id}{\mathrm{id}}
\renewcommand{\Re}{\operatorname{Re}}
\renewcommand{\Im}{\operatorname{Im}}
\newcommand{\rk}{\operatorname{rk}}
\newcommand{\rank}{\operatorname{rank}}
\newcommand{\vol}{\operatorname{vol}}
\newcommand{\const}{\mathrm{const}}
\newcommand{\half}{\tfrac{1}{2}}
\newcommand{\ie}{\textit{i.e.}}
\newcommand{\eg}{\textit{e.g.}}
\newcommand{\cf}{\textit{cf.}}
\newcommand{\from}{\leftarrow}
